% TOP_PROBLEM: {"id": 20001131, "title": "Compact Freiman models and relation peeling", "category_id": 2, "category": {"id": 2, "name": "combinatorics", "display_name": "Combinatorics", "description": "Counting problems, graph theory, discrete structures.", "slug": "combinatorics", "order_index": 2, "created_at": "2026-07-31T15:26:25.671Z"}, "status": "partially_solved", "problem_number": "AIM-COMBINATORICS-0256", "set_id": 14, "statusQualification": "Restores the exponent 2^{n-2}; length is diameter and isomorphism means a two-way Freiman isomorphism of order two."}
% FIRST_POSTED: 2026-10-01
% ENTRY_KIND: statement_only
% UnsolvedMath ID 20001131; source catalogue code AIM-COMBINATORICS-0256
% !TeX program = lualatex
% Definition and sources only; this file is not an LLM solution attempt.
\documentclass[11pt,a4paper]{article}
\usepackage[margin=25mm]{geometry}
\usepackage{fontspec}
\setmainfont{lmroman10-regular.otf}[BoldFont=lmroman10-bold.otf,
 ItalicFont=lmroman10-italic.otf,BoldItalicFont=lmroman10-bolditalic.otf]
\usepackage{amsmath,amssymb,mathtools}
\usepackage{unicode-math}
\setmathfont{latinmodern-math.otf}
\AtBeginDocument{\let\setminus\smallsetminus}
\usepackage{xurl,hyperref}
\hypersetup{colorlinks=true,urlcolor=blue}
\setcounter{secnumdepth}{0}
\setlength{\parindent}{0pt}
\setlength{\parskip}{5pt}
\setlength{\emergencystretch}{3em}
\begin{document}
\section*{Compact Freiman models and relation peeling}\label{20001131}
{\small UnsolvedMath ID 20001131 · AIM-COMBINATORICS-0256 · Combinatorics · AIM Workshop Problem Lists}
\subsection{Definitions and mathematical statement}
Let $A,B$ be finite subsets of $\mathbb Z$. A Freiman isomorphism
of order two is a bijection $\phi:A\to B$ such that, for every
$a_1,a_2,a_3,a_4\in A$, allowing repetitions,
\[
 a_1+a_2=a_3+a_4
 \quad\Longleftrightarrow\quad
 \phi(a_1)+\phi(a_2)=\phi(a_3)+\phi(a_4).
\]
Thus it preserves exactly all additive relations involving two
summands on each side, in both directions. The length of a nonempty
integer set $B$ means its diameter $\max B-\min B$, not its
cardinality or the number of integers in its containing interval.
Write $\ell(A)$ for the minimum length of an integer set Freiman
isomorphic to $A$.

The Konyagin--Lev conjecture in the workshop is that for every
integer $n\ge7$ and every $n$-element set $A\subseteq\mathbb Z$,
\[
                                  \ell(A)\le2^{n-2}.
\]
Equivalently, there must be such a bijection from $A$ to a subset
of the integer interval $[0,2^{n-2}]$; translation puts a minimal
model's smallest element at zero. The restriction $n\ge7$ is part
of the conjecture.
The exponential quantity arises from the model
$\{0,1,2,4,\ldots,2^{n-2}\}$, whose repeated doubling relations
motivate the proposed extremal length. A Freiman isomorphism is not
required to be the restriction of an affine map of $\mathbb Z$;
that extra requirement would change the problem. The exponent in
the interval endpoint is restored from the original PDF typography.
\subsection{Short English statement}
Can every n-element integer set, for n at least seven, be modeled with exactly the same additive pair relations inside an interval of length two to the power n minus two?
\subsection{Sources}
\begin{itemize}
\item[S1] ULAM.AI and the UnsolvedMath Contributors, supplied problems.json, record 20001131 (AIM-COMBINATORICS-0256), AIM Workshop Problem Lists. Catalogue identity and source transcription; CC BY 4.0.
\url{https://huggingface.co/datasets/ulamai/UnsolvedMath}.
\item[S2] American Institute of Mathematics, Recent trends in additive combinatorics, item 3.19. Original workshop question underlying AIM-COMBINATORICS-0256.
\url{https://aimath.org/WWN/additivecomb/additivecomb.pdf}.
\end{itemize}
\end{document}
